\clearpage\thispagestyle{empty}
\pdfbookmark[0]{Quellen zu Pallas. Komet. (17. August 1811)}{pallascomet1811aug-sources}
\begin{center}\textbf{Quellen zu Pallas. Komet. (17. August 1811)}\end{center}
\noindent Diese KI-Übertragung ins heutige Deutsch folgt dem Bericht in \href{https://gdz.sub.uni-goettingen.de/id/PPN236007467}{Gauß’ \textit{Werken}, Bd.~VI} (Göttingen, 1874), S.~327--330. Er erschien unter der Überschrift Göttingen in den \textit{Göttingischen gelehrten Anzeigen}, Stück 130 vom 17. August 1811, S.~\href{https://api.digitale-sammlungen.de/iiif/image/v2/bsb10538680_00265/full/1600,/0/default.jpg}{1289}, \href{https://api.digitale-sammlungen.de/iiif/image/v2/bsb10538680_00266/full/1600,/0/default.jpg}{1290}, \href{https://api.digitale-sammlungen.de/iiif/image/v2/bsb10538680_00267/full/1600,/0/default.jpg}{1291}, \href{https://api.digitale-sammlungen.de/iiif/image/v2/bsb10538680_00268/full/1600,/0/default.jpg}{1292} und \href{https://api.digitale-sammlungen.de/iiif/image/v2/bsb10538680_00269/full/1600,/0/default.jpg}{1293}. Das digitalisierte Exemplar stammt aus der Bayerischen Staatsbibliothek.
\par\smallskip\noindent Für den Unterschied der mittleren Länge im Jahr 1811 gibt die Werkausgabe auf S.~329 den Wert $-2.2$ an, der Erstdruck auf S.~1291 den Wert $-3{,}2$. Bei der fehlenden Deklination in der Berliner Beobachtung vom 18. Februar stehen in der Werkausgabe drei Striche, im Erstdruck drei Punkte. Bei den Elementen des Kometen steht in der Werkausgabe \textit{parobolischen}, im Erstdruck \textit{parabolischen}. KI-Übertragung und KI-Satz.
\clearpage
